\scanpage{151}
\mspage{9}
\deleted{tandis que le deuxième \unsure{mb|est} représente (compte tenu que
$\unsure{\mathbb R|\mathrm R}\mathrm{Hom}=\mathbb R\Gamma(C,\mathbb R\underline{\mathrm{Hom}})$}
\deleted{\[
 \mathrm{Hom}_{C}(G_C,R\underline{\mathrm{Hom}}(L\alpha_!(k_{C'}),F)
 \simeq\mathrm{Hom}_{C}(L\alpha_!(k_{C'}
\]}
L'isomorphie des $H^0$ des deux \unsure{mbs|mbres} est la formule d'adjonction
\[
 \mathrm{Hom}_{C'}(k_{C'},\alpha^*F)\simeq\mathrm{Hom}_C(\alpha_!(k_{C'}),F)
\]
d'où l'isomorphie pour les $H^i$ des deux mbs, par décalage.
De $=$ dans le cas de $\mathrm{HOT}$, l'isomorphie can. des $\pi_0$ des deux mbs est la formule d'adjonction.
Pour avoir l'isomorphisme plus fin dans $\mathbb D(e)$, la cat. \unsure{fondamentale} du dérivateur envisagé) en \unsure{termes|terme} de Hom \unsure{internes|interne},
on doit \deleted{\unsure{utili}} admettre l'existence d'objets unité (\inserted{les} $k_X$ dans un cas, \inserted{les} $e_X$ dans l'autre), compatible avec \unsure{image inverse|images inverses}, de telle façon
\[
 \mathrm{Hom\,int}(\unsure{\varepsilon|\mathcal E}_X,\underset{\substack{\rotatebox{-90}{$\in$}\\\mathbb D_X}}{F})\simeq F .
\]
On aura
\[
 \Gamma(C',\alpha^*F)\overset{\text{déf}}{\simeq} f'_*(\alpha^*(F))\simeq
 f_*\deleted{(\unsure{\mathrm{Hom\,int}}(\varepsilon_{C'},\unsure{\alpha^*}(F)))}\inserted{(\alpha_*\alpha^*F)} ,
\]
Et la question est donc de \unsure{donner|dégager} une interprétation de $\alpha_*\alpha^*F$ comme
\[
 \alpha_*\alpha^*F\simeq\mathrm{Hom\,int}(\alpha_!(\varepsilon_{C'}),F)
 \qquad\text{dans }\mathbb D(C)
\]
Pour ceci, on va écrire
\[
 \alpha^*F,\simeq\mathrm{Hom\,int}(\varepsilon_{C'},F),
\]
et la formule doit résulter d'une ``formule d'adjonction interne''
% archive sheet number: 151
